\section{The high estimate and its endpoint}\label{gex:high-section}

Let \(\beta_*\) be the supremum of \(1/2\) and the real parts of
nontrivial zeros of primitive finite-order Hecke \(L\)-functions over
\(F\). By \cite{OAI26}, \(\beta_*\le7/8\).
Suppose \(\beta_*>\theta\), and put
\(\Gamma=\beta_*-\theta>0\), \(\alpha=5/6\).

\subsection{Balancing the witness counts}

For rows \(q_u\asymp U=Z^d\), let \(R\) be the counting exponent:
\(\#\mathcal B\ll U^{R+\varepsilon}(1+T_1)^A\), with a fixed
height order \(A\). Use the detector bins of \cite[Section 8]{OAI26},
with real parameter \(a\) and freely chosen cutoff \(t\in[1,3/2]\).
In a nonfloor bin write \(\delta=2a-1\) and let
\(q\in[0,\delta/2]\) be the length-weighted prime amplitude; error
slots receive amplitude zero. The ceiling \(a\le\beta_*\) puts
actual bins in \(\delta\le3/4\); we retain the larger closed
rectangle for the endpoint certificate. We repeat the selection proof of
\cite[Proposition 19.2]{OAI26} with \(\widehat\kappa=3/4\).
Its printed hypothesis \(\beta_*>7/8\) is unnecessary for this
proof: the plain moment requires only
\(\beta_*\le(1+\widehat\kappa)/2\), already available here.

The coefficient convention matters. Put
\(\Theta=\langle\eta,\widehat{\mathcal T}\rangle\).
Conjugating entire witness and slot products when needed, the
prime coefficients are fixed combinations
\[
 \nu(p)\mathbf1_{p\in1_{\mathcal T}}
  =|\mathcal T|^{-1}
    \sum_{\vartheta\in\widehat{\mathcal T}}(\nu\vartheta)(p),
 \qquad \nu\in\Theta .
\]
Profiles and zero extensions are conjugated with their factors.
Large retained rows induce outside \(\Theta\), since they ramify
at a good prime; bounded unit rows are handled separately.

For witness lengths \(r<1\), \(m<1/2\), the inverse and doubled
plain capacities are \((1-r)/2\) and \(2(1-2m)/9\).
The witnesses satisfy
\[
 r+m\ge t-O(\varepsilon),\qquad
 r\le t+O(1/\log U)
\]
by \cite[Proposition 8.3]{OAI26}. Selecting whole slots in decreasing
amplitude gives the two count exponents
\[
 A_I(r)=1-\delta\{x+(1-x)r\},\qquad
 S_t(r)=1-\delta\{4x/9+(2-8x/9)(t-r)\},\quad x=q/\delta .
\]
The first decreases and the second increases with \(r\).
Their intersection gives
\[
 R_{\rm short}(t)=1-\delta+
       \frac{\delta P_x}{D_x}(3/2-t).
\]
For \(r\ge1\), inverse amplification and the upper bound for \(r\)
give \(L(t)=1-\delta+(\alpha-\delta)(t-1)\).
The unweighted plain bound handles \(m\ge1/2\).
Zero-capacity neighborhoods use these unweighted estimates,
not a marked estimate with vanishing decrement.
Thus the count is bounded by
\(\max\{R_{\rm short}(t),L(t)\}\), with the requested small losses.
The inequality \(r\le t+O(1/\log U)\) is essential here:
it bounds the increasing long-inverse exponent, rather than merely
asserting that one of the witnesses is large.

For \(1/50<\delta<\alpha\), balance these bounds by setting
\begin{equation}\label{gex:balanced-count}
\begin{gathered}
 y=1/2-x,\quad D_x=(37+34y)/18,\quad P_x=(7+18y+8y^2)/9,\\
 \mathcal J=(\alpha-\delta)D_x+\delta P_x
 =\{185+170y+(-138+12y+96y^2)\delta\}/108,\\
 t_*=1+\frac{\delta P_x}{2\mathcal J},\qquad
 R_*=1-\delta+\frac{(\alpha-\delta)\delta P_x}{2\mathcal J}.
\end{gathered}
\end{equation}
Indeed \(D_x(R_{\rm short}-L)=
\delta P_x/2-\mathcal J(t-1)\).
Since \(35/54\le\mathcal J\le5/2\), this yields uniformly
\begin{equation}\label{gex:row-count-with-loss}
 R\le R_*+c_{\rm row},
\end{equation}
with arbitrarily small prescribed witness, capacity and rounding
loss \(c_{\rm row}\). At \(\delta=\alpha\) use
\(R\le1-\delta+c_{\rm row}\); the floor \(a=51/100\) uses \(R=1\).

The new scales supply these moments: if
\(1/2\le d\le h+\zeta\) and \(0<\zeta<257/12000=5\ell-h\),
then \(\ell/d>1/5>7/37\).
Choose the moment mesh first, then sufficiently many even, disjoint
physical slots. The plain moment's auxiliary prime pool remains
separate. Their strict exponent gaps absorb annular
\(O(1/\log U)\) errors without changing any source mask.

\subsection{An exact endpoint inequality}

For sufficiently small \(e>0\), repeat the fixed-bin contour proof
of \cite[Lemma 20.1]{OAI26} on
\[
 (\Rea s,\Rea w,\Rea z)=(a+16e,1-a-6e,17/50).
\]
The global line \(\beta_*+20e\) and buffered rectangles remain
valid; \(a\ge51/100\) and \(\Rea(s+w)\ge1+10e\) admit the
unchanged first Euler region. Form amplitude and witness sets
only on the retained contour. Joint conductor allocation bounds
the numerator and all error slots together.

The actual Mellin and slot powers give, relative to \(C(\theta)\),
\begin{equation}\label{gex:high-central-exponent}
 E_\theta(d)=a-\theta+h(17/50-1/6)-a l_y-\ell/2+q\ell
             +d(R+\delta/2-17/50).
\end{equation}
In particular \(q\ell\) is already measured at base \(Z\).

\begin{lemma}\label{gex:endpoint}
For \(0\le\delta\le5/6\), \(0\le x\le1/2\), and \(p=1/6000\),
substitution of \(q=\delta x\), \(R=R_*\) gives
\(E_\theta(h)\le-1/200000\).
\end{lemma}
\begin{proof}
Write \(E_0=-1/48+(2/3)\delta+\delta x/6-(13/16)(1-R_*)\).
The new geometry gives exactly
\[
 E_\theta(h)=E_0+p\{-1/4+\delta(1+x)+(3/2)R_*\}.
\]
For \(m=1/200000\), multiplication gives
\(10368\mathcal J(-E_\theta(h)-m)=A\delta^2+B\delta+C\), where
\begin{align*}
 A&=6(32032y^3+94028y^2+131008y+51021)/125,\\
 B&=-(650644y^2+9423268y+5918793)/3125,\\
 C&=6186(34y+37)/625,
\end{align*}
and
\[
 4AC-B^2=\frac{\begin{aligned}
 &19787957489264y^4+69061294318816y^3\\[-2pt]
 &+50726295758792y^2+10600196228952y+1254989151
 \end{aligned}}{9765625}>0 .
\]
Since \(A,\mathcal J>0\), completing the square proves the claim:
\(4A(A\delta^2+B\delta+C)=(2A\delta+B)^2+4AC-B^2>0\).
\end{proof}

\subsection{All moderate ranges}

The count envelopes have \(1-\delta\le R\le3/2\); their frequency
slope \(R+\delta/2-17/50\) lies between \(4/25\) and \(2\).
Thus \(d=h\) controls \(1/2\le d\le h\); extending to \(h+\zeta\)
costs \(2\zeta\). The other branches give
\[
\begin{array}{ll}
 \text{floor:}&E_\theta(h)\le-7/1200+15p/8,\\
 \delta=5/6,\ R=1-\delta:
   &E_\theta(h)\le-1/48-\delta/16+15p/8,\\
 1/100\le d\le1/2:
   &E_\theta(d)\le-49/14400+139p/150 .
\end{array}
\]
The last uses \(R=76/75-2\delta/3\), whose slope is positive
and whose geometry derivative is \(a/2+1/100+q\le139/150\).
All three margins exceed \(m\).

If nominal row and frequency exponents \(d,v\) differ, with
\(v\le h+\zeta\), \(d-v\le\mu\), their extra cost is at most
\(2\zeta+(3/2)\mu\). Including contour shifts, amplitude width
\(\vartheta\), central power allowance \(\varepsilon_0\), and
the remaining finite loss \(\lambda_{\rm loss}\), the total is bounded by
\begin{equation}\label{gex:central-loss-budget}
 \Lambda=h c_{\rm row}+2\zeta+(3/2)\mu+
          26e+(K+8)\varepsilon_0+\vartheta/5+\lambda_{\rm loss} .
\end{equation}
Here \(26\) bounds \(16-6l_y+12v\), and \(\ell<1/5\).
The compact parameter ranges and positive denominator bounds make
every loss uniform. Choose their budgets before the target,
the mesh before \(K\), and \(\varepsilon_0\) after \(K\), so that
\(\Lambda<m/4\). Reserving another \(m/4\) for the normalizer
and finite dyadic multiplicities leaves a positive saving.
Comparison with \(C(\beta_*)\) subtracts the further positive gap
\(\Gamma\). This accounts for every moderate nonprincipal row.
